\chapter{Conclusion}
\label{chap:conclusion}

The work of this thesis focused on developing a deep learning model for malaria prescreening with limited-resource conditions in mind, including the future use of smartphone-based images. This thesis is part of a larger system that also includes the development of an ocular adapter, which is intended to be attached to a microscope and enable the capture of thick blood smear images using a smartphone. The model developed in this work is intended to classify complete images as infected or uninfected. An important requirement was therefore also the evaluation of lightweight architectures that could later be suitable for deployment.

To achieve the intended goal of this work, the NLM/NIH thick blood smear dataset was selected for training, validation and testing. A patient-based split of 70\% training, 15\% validation and 15\% testing was used. The selected architectures were ResNet50 as the main high-performance reference model, ResNet18 as a smaller ResNet architecture, and MobileNetV3-Small as the deployment-relevant lightweight model. Since the different parts of the overall project were developed in parallel, the final image resolution of the ocular adapter was not yet available during this thesis. For this reason, different resolution levels were simulated, starting from HR at 224 pixels and reducing the image resolution to LR3 at 56 pixels before resizing back to the model input size. This allowed the behaviour of the models under increasingly reduced image resolution to be investigated. Grad-CAM and Grad-CAM++ were also used to analyse which image regions influenced the predictions. In addition to the NLM/NIH test data, an external test using the AIR Uganda dataset was performed to evaluate the models on images with different characteristics from the development dataset.

The testing results showed that all architectures achieved very high performance at full resolution. ResNet50 reached an accuracy of 99.58\% at HR. Its performance remained high at LR1 and LR2, but a clearer decrease was observed at LR3, where the accuracy dropped to 92.86\%. The specificity also decreased to 73.89\%, while the sensitivity remained high at 98.00\%. ResNet18 and MobileNetV3-Small both reached the highest full-resolution accuracy of 99.69\% and a sensitivity of 100\%, meaning that all infected images in the internal test set were correctly classified. Under stronger resolution reduction, the main decrease was seen especially in the classification of uninfected images, with more false-positive predictions being produced. MobileNetV3-Small was therefore further optimised at LR3 and finally reached an accuracy of 92.76\%, a sensitivity of 96.40\%, a specificity of 79.31\% and an F1-score of 95.45\%.

The Grad-CAM and Grad-CAM++ results raised more questions about the behaviour of the models. In several examples, strong activations were visible around the edges and dark background of the microscope field instead of only inside the blood smear region. This raised the possibility that the models may partly use image characteristics that are not directly related to the parasites. To further investigate this, the external AIR Uganda dataset was used because its images have different visual characteristics compared with the NLM/NIH dataset. Since this dataset contains only infected images, only the sensitivity could be evaluated. ResNet50 reached a sensitivity of 100\% at HR and 98.48\% at LR3, with only 18 infected images missed at LR3. MobileNetV3-Small reached a sensitivity of 95.60\% at LR3, while ResNet18 showed a lower sensitivity of 82.32\% with 209 false-negative predictions. These results support that the models are not relying only on the specific borders or background characteristics of the NLM/NIH images. However, shortcut learning cannot be completely excluded. The external test also cannot show whether the models tend to predict the infected class when uncertain, because no uninfected samples were available in the AIR Uganda dataset.

Overall, the whole-image classification approach is technically promising based on the results of this work and the comparison with previous studies. Even when the image resolution was strongly reduced, good performance was still achieved. MobileNetV3-Small was the most relevant architecture for future deployment and also showed good performance at both high and low resolution. 

The work of this thesis should be understood as a research prototype for future development and is not yet clinically validated for real-life malaria prescreening. The next important step is to obtain images captured with the combination of the developed ocular adapter, microscope and smartphone. These images will provide the real image resolution and characteristics that can be expected during future deployment. The models can then be tested using real captured images instead of only simulated resolution reductions. Future testing should also include more external datasets, especially datasets containing uninfected images, since this remained an open question during the testing phase. This would allow the specificity and false-positive behaviour of the models to be evaluated more reliably. In addition, using more uninfected samples during future model development could potentially help reduce the number of uninfected images that are incorrectly classified as infected.

A possible future extension of the system could also be parasite detection. After the binary prescreening approach has been further evaluated and improved, the system could be extended to identify parasites inside the image and later also distinguish between different Plasmodium species.
